%% ma: L12-B14-0125
%% lop: 12
%% chuong: 5
%% bai: 14
%% dang: 12.14.35
%% loai: TLN
%% muc: VD
%% dap_an: "36"
%% nguon: "Ảnh giáo viên gửi 07/11/2026 (Tư Duy Mở - tuduymo.com, đề tự luyện Phương trình mặt phẳng, đề khoảng cách), Câu 21"
%% kien_thuc: "Bài 14 (phương trình mặt phẳng qua một điểm, biết vectơ pháp tuyến; khoảng cách từ điểm đến mặt phẳng), quan hệ đường vuông góc và đường xiên (lớp 11)"
%% trang_thai: cho-duyet
%% ly_do: "Dạng 12.14.35 mới đề xuất, chờ giáo viên duyệt."
%% ghi_chu: "Tài liệu gốc không kèm đáp án; đáp án do Claude giải và kiểm bằng sympy."
\begin{ex}
Trong không gian $Oxyz$, cho hai điểm $A(3;1;0)$, $B(-2;0;3)$. Mặt phẳng $(P)$ đi qua $A$ và khoảng cách từ $B$ đến $(P)$ lớn nhất có phương trình tổng quát tương ứng là $(P)\colon ax+2y+cz+d=0$. Hãy xác định giá trị của biểu thức $T=a+c-d$.
\shortans{36}
\loigiai{Gọi $H$ là hình chiếu của $B$ trên $(P)$: $d\left(B,(P)\right)=BH\le BA$, dấu bằng khi $(P)\perp AB$ tại $A$.
$\overrightarrow{AB}=(-5;-1;3)$, chọn vectơ pháp tuyến $(10;2;-6)$: $(P)\colon 10(x-3)+2(y-1)-6z=0\Leftrightarrow10x+2y-6z-32=0$.
$T=10-6+32=36$.}
\end{ex}
